\section{Introduction}
\label{sec:intro}

\PARstart{A}{utomatic} estimation of student engagement and attention has been studied for more than a decade, first from facial expressions~\cite{whitehill2014faces} and then from webcam recordings of individual learners such as DAiSEE~\cite{gupta2016daisee} and EngageNet~\cite{singh2023engagenet}. Recent systems combine several modalities (the face, facial landmarks, head pose and, when available, physiology) through attention or gating~\cite{baltrusaitis2019multimodal,arevalo2017gated,tran2026ordinal}, and data fusion is a recurrent theme of multimodal learning analytics~\cite{chango2022datafusionmmla}. The underlying premise is that each stream contributes information and that a learned fusion weights the streams appropriately.

This premise is rarely tested in isolation. Reported gains are usually accuracy differences on a single split, obtained with different backbones and protocols for the fused and the single-stream systems. The multimodal learning literature shows that a jointly trained model can underperform the best separately trained unimodal model, because the added capacity overfits and the modalities overfit and generalise at different rates~\cite{wang2020makes}, and because a dominant modality suppresses the optimisation of the others~\cite{peng2022balanced}. Engagement studies contain similar observations: the EngageNet baselines obtain a lower test accuracy with fused features than with one feature set~\cite{singh2023engagenet}, a study in secondary-school classrooms finds score-level fusion on par with the best modality~\cite{sumer2023multimodal}, and a recent transformer study on classroom data finds its fused model within the variability of its visual model~\cite{fragkiadakis2026fairness}. Whether fusion helps for classroom attention, where labels are noisy, ordinal and imbalanced, has nevertheless not been measured under a single controlled protocol.

A second, related question concerns the evaluation itself. Scores depend on what the test set shares with the training set, and such dependence is a documented cause of optimistic results in many fields~\cite{kapoor2023leakage}. In clinical and wearable sensing, splitting records rather than subjects is known to inflate performance~\cite{saeb2017need}, as does placing neighbouring windows on both sides of a split~\cite{hammerla2015lets}, and model-selection bias can be as large as the difference between learning algorithms~\cite{cawley2010overfitting}. Classroom data add a further level of structure: students of the same class attend the same lesson at the same time, so that windows of different students are not independent either.

DIPSER~\cite{marquezcarpintero2025dipser} makes it possible to study both questions. It provides in-person classroom sessions with a camera and a smartwatch per student, and attention ratings from several annotators. Because a controlled study requires hundreds of trainings, we encode every stream with discrete Tchebichef moments~\cite{mukundan2001tchebichef}, which require no training and keep each downstream model small. The study therefore runs without a GPU cluster, and every configuration can be repeated over folds and seeds.

This article reports a negative result and its analysis; it does not propose a new state of the art. It addresses six questions: (Q1) whether any fusion exceeds the best single stream; (Q2) whether a learnable gate differs from mean pooling; (Q3) whether a linear input-to-head link helps; (Q4) which ordinal objective performs best under rater disagreement; (Q5) whether the answers transfer to DAiSEE and EngageNet; and (Q6) whether they hold on the full DIPSER release, and how much of a reported score depends on the evaluation protocol.

The contributions are the following.
\begin{enumerate}
\item A controlled fusion study on DIPSER with one instrumented model, nested epoch selection, three seeds and a permuted-label control (495 fits), in which no fusion configuration exceeds the best single stream and the learned gate does not differ from mean pooling. The same comparison is repeated on DAiSEE and EngageNet.
\item A replication on the full DIPSER release (\fcPersons{} students) with held-out students, followed by seven attempts to improve the sensor model. Apart from one exploratory search, the success criterion of each attempt was fixed before the run, and none succeeds.
\item A measurement of the effect of the evaluation split with the model held fixed (random windows, held-out students, held-out class), together with label-only references and a label-persistence baseline that put the sensor scores in context.
\item The complete record of the study: per-seed results, out-of-fold predictions, the analysis plan with its criteria, and scripts that regenerate every table and figure.
\end{enumerate}

The remainder of the article is organised as follows. Section~\ref{sec:related} reviews related work. Section~\ref{sec:data} describes the datasets and the task, Section~\ref{sec:instrument} the model and Section~\ref{sec:protocol} the protocol. Section~\ref{sec:results} reports the results on the DIPSER slice and on the two webcam datasets, and Section~\ref{sec:fullcorpus} the results on the full release. Section~\ref{sec:discussion} discusses the findings and their limitations, and Section~\ref{sec:conclusion} concludes.
